\section{Conclusion and Future Work}
\label{sec:conclusion}

We report Stage 2b of a pre-registered pilot study of LLM-agent Bertrand pricing dynamics, with a headline finding of comparable terminal profit gain ($\Delta_{\text{profit}} = 0.4932$ at $n = 5$ versus $0.4522$ at $n = 2$, both capturing approximately half the Nash-to-Monopoly profit gap) coupled to sharply divergent convergence rates ($30/30$ at $n = 5$ versus $1/15$ at $n = 2$ within a 50-round horizon). The asymmetry is consistent with a wider basin of attraction at smaller agent counts and is the central empirical contribution of this paper.

The methodological contribution is the layered integrity stack: cryptographic pre-registration on OSF Registries with RFC 3161 trusted timestamps as the foundation, independent cross-toolchain reproduction (TypeScript primary analyzer plus Python/SciPy reference verifier) as a procedural reinforcing layer, and structured multi-dimensional internal review with adversarial calibration as an iterative reinforcing layer. The three commitments are not equivalent: cryptographic pre-registration is the strongest integrity primitive available for analysis-plan locking; the verification and adversarial layers reinforce it but cannot substitute for it. The stack caught one analyzer bug pre-submission (a Fisher's exact denominator error that would have inverted the host-effect verdict); fired one pre-committed downgrade branch literally on the data ($k = 4$ to $k = 2$ at GATE-5); and produced six required revisions to the analysis prior to this preprint, three of which materially tightened the published claim wording. The protocol is replicable: every layer is open-source and the procedure generalizes to any pre-registered multi-agent LLM study.

Stage 3 expands $n_{\text{rep}}$ to $\geq 30$ for both conditions and lengthens the horizon to 100 rounds, and instruments the trace schema with \texttt{model\_version}, \texttt{region}, \texttt{ttft\_ms}, \texttt{tpot\_ms}, \texttt{cache\_hit\_tokens}, \texttt{cache\_miss\_tokens}, and \texttt{auth\_handshake\_ms} per Addendum \#1 \S F to isolate LLM-serving variability from reasoning-depth variability. Three falsification conditions are pre-locked in this paper (\S\ref{sec:discussion-falsifiability}): the basin-width-by-agent-count hypothesis, the $\Delta_{\text{profit}}$ asymmetry magnitude, and the LLM-versus-$Q$-learner family-generalization claim. Stage 4 conducts cross-model replication across Haiku sizes, Sonnet, Opus, and non-Anthropic frontier models to formally test the simplicity-bias mechanism (\citealp{wu2025whenMoreIsLess}; \S\ref{sec:discussion-wu2025}) and the family-generalization claim.

The codebase, the pre-registered analysis pipeline, the four RFC 3161-stamped addenda, and the full Verification \& Correction Log are publicly available at \url{https://github.com/HHA-Applied-Research-Institute/colludebench-cascade}. We invite Stage 3 replication --- particularly at the duopoly limit where the Stage 2b power envelope is thinnest --- and the parallel $Q$-learner positive-control experiment pre-committed in \texttt{pilot/VALIDATION\_PROTOCOL.md} as the cross-family discriminator that this pilot defers to.
